% =========================
% L23.tex
% =========================

\chapter{L23 -- Principio di invarianza di LaSalle, esercizi ed elementi di sintesi con Lyapunov}
\label{chap:L23}

\section{Insiemi invarianti}

Prima di enunciare il principio di LaSalle conviene richiamare con precisione il concetto di insieme invariante.

\medskip

\noindent
\textbf{Definizione.}
Sia dato un sistema dinamico autonomo.

\begin{itemize}
    \item Nel caso continuo:
    \[
    \dot x = f(x), \qquad x\in\mathbb{R}^n.
    \]
    \item Nel caso discreto:
    \[
    x(k+1)=f(x(k)).
    \]
\end{itemize}

Un insieme \(M\subseteq \mathbb{R}^n\) si dice \textbf{positivamente invariante} se ogni traiettoria che parte in \(M\) rimane in \(M\) per tutti i tempi futuri.

Nel caso continuo questo significa:
\[
x(t_0)\in M
\quad \Longrightarrow \quad
\varphi(t,t_0,x(t_0))\in M
\qquad \forall t\ge t_0.
\]

Nel caso discreto:
\[
x(0)\in M
\quad \Longrightarrow \quad
x(k)\in M
\qquad \forall k\ge 0.
\]

\medskip

\noindent
In modo intuitivo: un insieme è invariante se, una volta entrati dentro, non se ne esce più.

\begin{figure}[h!]
\centering
\begin{tikzpicture}[scale=1]
    \draw[thick] (0,0) ellipse (3 and 2);
    \node at (0,2.35) {$M$};

    \fill (-1,0.3) circle (1.5pt);
    \fill (1.2,-0.5) circle (1.5pt);

    \draw[-{Latex[length=3mm]},thick] (-1,0.3) .. controls (-0.2,0.8) and (0.4,0.6) .. (1.2,-0.5);
    \draw[-{Latex[length=3mm]},thick] (1.2,-0.5) .. controls (1.8,-0.3) and (2.0,0.4) .. (1.4,1.0);

    \draw[-{Latex[length=3mm]},dashed] (-2.7,-0.7) -- (-2.2,-0.1);
    \node at (-3.2,-0.9) {\small traiettoria};
\end{tikzpicture}
\caption{Schema qualitativo: una traiettoria che parte in \(M\) rimane in \(M\).}
\end{figure}

\section{Perché serve LaSalle}

Il teorema diretto di Lyapunov classico dice:

\begin{itemize}
    \item se esiste \(V\) positiva definita e \(L_fV\le 0\), allora l'equilibrio è stabile;
    \item se esiste \(V\) positiva definita e \(L_fV<0\) per ogni \(x\neq 0\), allora l'equilibrio è asintoticamente stabile.
\end{itemize}

Il punto delicato è il caso intermedio:
\[
L_fV(x)\le 0
\quad \text{ma non strettamente } <0.
\]
In questa situazione il teorema diretto, da solo, spesso garantisce soltanto la stabilità, non la convergenza all'equilibrio.

Il principio di invarianza di LaSalle serve esattamente a colmare questo vuoto: ci dice \emph{dove vanno a finire}, nel lungo periodo, le traiettorie di un sistema quando la funzione di Lyapunov non cresce ma può rimanere costante su qualche sottoinsieme.

\section{Enunciato del principio di invarianza di LaSalle}

\subsection{Versione per sistemi a tempo continuo}

Consideriamo il sistema
\[
\dot x = f(x), \qquad x\in\mathbb{R}^n,
\]
con \(f\) sufficientemente regolare da garantire esistenza e unicità delle soluzioni.

Sia \(\Omega\subseteq \mathbb{R}^n\) un insieme chiuso, limitato e positivamente invariante.  
Sia inoltre
\[
V:\Omega\to\mathbb{R}
\]
una funzione di classe \(C^1\) tale che
\[
L_fV(x)=\nabla V(x)^\top f(x)\le 0
\qquad \forall x\in\Omega.
\]

Definiamo l'insieme
\[
E=\{x\in\Omega : L_fV(x)=0\}.
\]
Sia poi \(M\) il \textbf{massimo insieme invariante} contenuto in \(E\).

\medskip

\noindent
\textbf{Teorema di LaSalle.}
Per ogni traiettoria che parte in \(\Omega\),
\[
x(0)\in\Omega,
\]
si ha
\[
\operatorname{dist}(x(t),M)\to 0
\qquad \text{per } t\to+\infty.
\]

\medskip

\noindent
In parole: la traiettoria converge al massimo insieme invariante contenuto nella zona in cui \(V\) smette di diminuire.

\subsection{Versione per sistemi a tempo discreto}

Per un sistema discreto
\[
x(k+1)=f(x(k)),
\]
si ragiona in modo del tutto analogo.  
Se esiste una funzione \(V\) tale che
\[
\Delta V(x)=V(f(x))-V(x)\le 0,
\]
allora si considera l'insieme
\[
E=\{x\in\Omega : \Delta V(x)=0\},
\]
e il massimo insieme invariante \(M\subseteq E\).  
Anche in questo caso le traiettorie convergono a \(M\).

\section{Interpretazione geometrica}

Il principio di LaSalle si capisce bene così.

\begin{itemize}
    \item La funzione \(V(x)\) si comporta come una ``energia''.
    \item Se \(L_fV\le 0\), questa energia non aumenta.
    \item Quindi la traiettoria non può vagare in modo arbitrario: è costretta a scendere, o al massimo a restare sullo stesso livello di energia.
    \item Nel lungo periodo, allora, il sistema deve finire nella regione in cui la discesa si arresta, cioè in
    \[
    E=\{x:L_fV(x)=0\}.
    \]
    \item Ma non tutto \(E\) è necessariamente compatibile con la dinamica: la traiettoria può fermarsi solo nel \emph{massimo insieme invariante} contenuto in \(E\).
\end{itemize}

\begin{figure}[h!]
\centering
\begin{tikzpicture}[scale=1]
    \draw[->] (-3.2,0) -- (3.2,0) node[right] {$x_1$};
    \draw[->] (0,-3.2) -- (0,3.2) node[above] {$x_2$};

    \draw[thick] (0,0) circle (2.4);
    \draw[thick] (0,0) circle (1.7);
    \draw[thick] (0,0) circle (1.0);

    \draw[red,thick] (-2.8,0.8) .. controls (-1.6,1.1) and (-0.5,0.8) .. (0.1,0.2);
    \draw[-{Latex[length=3mm]},red,thick] (-0.05,0.25) -- (0.1,0.2);

    \draw[blue,thick] (-1.3,-1.3) -- (1.3,1.3);
    \node[blue] at (1.7,1.6) {\small \(E=\{L_fV=0\}\)};

    \fill (0,0) circle (1.7pt);
    \node at (0.35,-0.3) {\small \(M\)};
\end{tikzpicture}
\caption{Idea qualitativa: la traiettoria scende lungo le curve di livello di \(V\) e converge al massimo insieme invariante contenuto in \(E\).}
\end{figure}

\section{Corollario fondamentale per la stabilità asintotica}

Il caso più importante, in pratica, è questo.

\medskip

\noindent
\textbf{Corollario.}
Se l'origine è un equilibrio, esiste una funzione \(V\) positiva definita e
\[
L_fV(x)\le 0,
\]
e inoltre il massimo insieme invariante contenuto in
\[
E=\{x:L_fV(x)=0\}
\]
è il solo punto
\[
M=\{0\},
\]
allora l'origine è \textbf{asintoticamente stabile}.

\medskip

\noindent
Questo risultato è estremamente utile perché permette di concludere l'asintotica stabilità anche quando \(L_fV\) è solo semidefinita negativa.

\section{Mini-esempio: convergenza a una retta di equilibri}

Consideriamo il sistema
\[
\begin{cases}
\dot x_1 = -x_1,\\[4pt]
\dot x_2 = 0.
\end{cases}
\]
Scegliamo
\[
V(x)=\frac{1}{2}x_1^2.
\]
Allora
\[
L_fV=x_1\dot x_1=-x_1^2\le 0.
\]
Quindi
\[
E=\{x:L_fV=0\}=\{x_1=0\}.
\]
Ma l'insieme \(\{x_1=0\}\) è esso stesso invariante, perché se \(x_1(0)=0\), allora \(x_1(t)=0\) per ogni \(t\), mentre \(x_2\) resta costante.

Dunque il massimo insieme invariante contenuto in \(E\) è
\[
M=\{(0,x_2):x_2\in\mathbb{R}\}.
\]
Per LaSalle, le traiettorie convergono a questa retta, non necessariamente all'origine.

Questo esempio è importante perché mostra bene che:
\begin{center}
LaSalle non dice automaticamente che si converge a un punto; dice che si converge al massimo insieme invariante contenuto in \(E\).
\end{center}

\section{Collegamento con l'esercizio della lezione precedente}

Riprendiamo il sistema già studiato in precedenza, nel caso \(k=0\):
\[
\begin{cases}
\dot z_1 = -z_1-z_2^2,\\[4pt]
\dot z_2 = 4z_1z_2.
\end{cases}
\]
Scegliamo
\[
V(z)=\frac{1}{2}\left(4z_1^2+z_2^2\right).
\]
Allora
\[
L_fV = -4z_1^2\le 0.
\]
Quindi
\[
E=\{z:L_fV(z)=0\}=\{z_1=0\}.
\]

Per capire il massimo insieme invariante contenuto in \(E\), imponiamo \(z_1(t)=0\) per tutti i tempi.  
Ma allora dalla prima equazione segue
\[
\dot z_1 = -z_2^2.
\]
Affinché \(z_1\) resti identicamente nullo, deve valere necessariamente
\[
-z_2^2=0 \quad \Longrightarrow \quad z_2=0.
\]
Quindi l'unico insieme invariante contenuto in \(E\) è il solo punto
\[
M=\{(0,0)\}.
\]
Per il corollario di LaSalle segue allora che l'origine è \textbf{asintoticamente stabile}.

Questo è esattamente il tipo di situazione in cui LaSalle è più potente del teorema diretto classico.

\section{Esercizio 1: studio di un sistema non lineare}
\label{sec:L23-es1}

Studiamo il sistema
\begin{equation}
\label{eq:L23-es1}
\begin{cases}
\dot x_1 = -x_1^3 + x_2,\\[4pt]
\dot x_2 = -x_1^3 - x_2^3.
\end{cases}
\end{equation}

\subsection{Equilibri}

Gli equilibri si trovano imponendo
\[
\begin{cases}
0 = -x_1^3 + x_2,\\[4pt]
0 = -x_1^3 - x_2^3.
\end{cases}
\]
Dalla prima equazione:
\[
x_2=x_1^3.
\]
Sostituendo nella seconda:
\[
0=-x_1^3-(x_1^3)^3=-x_1^3-x_1^9=-x_1^3(1+x_1^6).
\]
Poiché
\[
1+x_1^6>0 \qquad \forall x_1\in\mathbb{R},
\]
segue
\[
x_1=0.
\]
Di conseguenza, anche
\[
x_2=0.
\]
L'unico equilibrio è quindi
\[
(0,0).
\]

\subsection{Metodo indiretto}

Il Jacobiano è
\[
A(x)=
\begin{pmatrix}
-3x_1^2 & 1\\[4pt]
-3x_1^2 & -3x_2^2
\end{pmatrix}.
\]
Valutando nell'origine:
\[
A(0,0)=
\begin{pmatrix}
0 & 1\\
0 & 0
\end{pmatrix}.
\]
Il polinomio caratteristico è
\[
\lambda^2=0,
\]
quindi gli autovalori sono entrambi nulli. Il metodo indiretto è dunque inconcludente.

\subsection{Scelta di una funzione di Lyapunov}

Una scelta quadratica standard,
\[
V=\frac{1}{2}(x_1^2+x_2^2),
\]
non è conveniente perché genera termini misti difficili da controllare:
\[
L_fV=x_1(-x_1^3+x_2)+x_2(-x_1^3-x_2^3).
\]
Per eliminare i termini misti e adattarsi alle potenze del sistema conviene scegliere una funzione di Lyapunov \emph{non quadratica}.

Una scelta efficace è
\[
V(x)=\frac{1}{4}x_1^4+\frac{1}{2}x_2^2.
\]
Essa è positiva definita e radialmente illimitata.

Calcoliamo la derivata lungo le traiettorie:
\[
L_fV = x_1^3\dot x_1 + x_2\dot x_2.
\]
Sostituendo \eqref{eq:L23-es1}:
\begin{align*}
L_fV
&= x_1^3(-x_1^3+x_2)+x_2(-x_1^3-x_2^3)\\
&= -x_1^6 + x_1^3x_2 - x_1^3x_2 - x_2^4\\
&= -x_1^6 - x_2^4.
\end{align*}
Quindi
\[
L_fV(x)<0
\qquad \forall x\neq 0.
\]
Essendo \(V\) positiva definita e radialmente illimitata, e \(L_fV\) negativa definita, concludiamo che l'origine è \textbf{globalmente asintoticamente stabile}.

\medskip

\noindent
\textbf{Osservazione metodologica.}
Questo esercizio è istruttivo perché mostra che, nei sistemi non lineari, la funzione di Lyapunov non deve per forza essere quadratica: spesso bisogna scegliere le potenze giuste per adattarsi alla struttura del campo vettoriale.

\section{Esercizio 2: pendolo smorzato}
\label{sec:L23-pendolo}

Consideriamo il pendolo smorzato
\begin{equation}
\label{eq:L23-pendolo}
\begin{cases}
\dot x_1 = x_2,\\[4pt]
\dot x_2 = -\dfrac{c}{m\ell^2}x_2 - \dfrac{g}{\ell}\sin x_1,
\end{cases}
\end{equation}
dove:
\begin{itemize}
    \item \(x_1=\theta\) è l'angolo,
    \item \(x_2=\dot\theta\) è la velocità angolare,
    \item \(m\) è la massa,
    \item \(\ell\) la lunghezza,
    \item \(c\ge 0\) il coefficiente di smorzamento.
\end{itemize}

\subsection{Linearizzazione nell'origine}

L'origine \((0,0)\) è un equilibrio perché
\[
x_2=0, \qquad \sin x_1 = 0 \text{ per } x_1=0.
\]
Il Jacobiano nell'origine è
\[
A(0,0)=
\begin{pmatrix}
0 & 1\\[4pt]
-\dfrac{g}{\ell} & -\dfrac{c}{m\ell^2}
\end{pmatrix}.
\]

\begin{itemize}
    \item Se \(c>0\), il sistema linearizzato è asintoticamente stabile.
    \item Se \(c=0\), il linearizzato ha autovalori puramente immaginari e il metodo indiretto non conclude.
\end{itemize}

\subsection{Funzione energia}

La funzione naturale da provare è l'energia meccanica:
\[
V(x)=mg\ell\bigl(1-\cos x_1\bigr)+\frac{m\ell^2}{2}x_2^2.
\]
Vicino all'origine questa funzione è positiva definita, perché:
\[
1-\cos x_1 \sim \frac{x_1^2}{2}
\qquad \text{per } x_1\to 0.
\]

La derivata lungo le traiettorie è
\[
L_fV
=
mg\ell\sin x_1\,\dot x_1 + m\ell^2 x_2\dot x_2.
\]
Sostituendo \eqref{eq:L23-pendolo}:
\begin{align*}
L_fV
&=
mg\ell\sin x_1\,x_2
+
m\ell^2x_2\left(
-\frac{c}{m\ell^2}x_2 - \frac{g}{\ell}\sin x_1
\right)\\
&=
mg\ell\sin x_1\,x_2
-cx_2^2
-mg\ell x_2\sin x_1\\
&=
-cx_2^2.
\end{align*}
Quindi:
\[
L_fV = -cx_2^2 \le 0.
\]

\subsection{Caso \(c>0\): applicazione di LaSalle}

Se \(c>0\), allora
\[
L_fV=0 \iff x_2=0.
\]
Quindi
\[
E=\{(x_1,x_2):x_2=0\}.
\]

Per applicare bene LaSalle conviene limitarsi a un sottolivello di energia che non includa il punto alto del pendolo, cioè scegliere
\[
\Omega_\rho = \{x: V(x)\le \rho\}
\qquad \text{con } 0<\rho<2mg\ell.
\]
In questo modo restiamo in una regione attorno all'origine e non includiamo l'equilibrio instabile in \(x_1=\pi\).

Se una traiettoria resta in \(E\), allora \(x_2(t)=0\) per ogni \(t\). Ma dalla seconda equazione di \eqref{eq:L23-pendolo} segue
\[
\dot x_2 = -\frac{g}{\ell}\sin x_1.
\]
Affinché \(x_2\) resti identicamente nullo, deve essere
\[
\sin x_1 = 0.
\]
Nel sottolivello scelto, l'unico punto che soddisfa questa condizione è
\[
x_1=0.
\]
Dunque il massimo insieme invariante contenuto in \(E\cap \Omega_\rho\) è il solo punto
\[
M=\{(0,0)\}.
\]
Per LaSalle, l'origine è \textbf{asintoticamente stabile}.

\subsection{Caso \(c=0\): solo stabilità}

Se \(c=0\), allora
\[
L_fV=0
\qquad \forall x.
\]
Quindi il teorema di LaSalle non può dare convergenza all'origine, perché tutto l'insieme considerato appartiene a \(E\). In effetti il sistema è conservativo: l'energia si conserva e le traiettorie, in generale, non convergono all'equilibrio.

Quindi:
\[
c=0
\quad \Longrightarrow \quad
(0,0) \text{ è stabile, ma non asintoticamente stabile.}
\]

\section{Esercizio 3: sistema discreto}
\label{sec:L23-TD}

Consideriamo il sistema discreto
\begin{equation}
\label{eq:L23-TD}
\begin{cases}
x_1(k+1)=\alpha x_2(k)-x_1^2(k)x_2(k),\\[4pt]
x_2(k+1)=-\alpha x_1(k)+x_1(k)x_2^2(k).
\end{cases}
\end{equation}

\subsection{Linearizzazione nell'origine}

L'origine è un equilibrio.  
Il linearizzato nell'origine è
\[
x(k+1)=Ax(k),
\qquad
A=
\begin{pmatrix}
0 & \alpha\\
-\alpha & 0
\end{pmatrix}.
\]
Gli autovalori sono
\[
\lambda_{1,2}=\pm j\alpha,
\]
quindi il loro modulo è
\[
|\lambda_{1,2}|=|\alpha|.
\]

Da qui segue:
\[
|\alpha|<1 \Rightarrow (0,0)\text{ asintoticamente stabile},
\]
\[
|\alpha|>1 \Rightarrow (0,0)\text{ instabile},
\]
\[
|\alpha|=1 \Rightarrow \text{metodo indiretto inconcludente}.
\]

\subsection{Caso critico \(\alpha=1\)}

Poniamo ora \(\alpha=1\).  
Proviamo con
\[
V(x)=x_1^2+x_2^2.
\]
Calcoliamo
\[
\Delta V(x)=V(x(k+1))-V(x(k)).
\]
Sostituendo la dinamica:
\begin{align*}
\Delta V
&=
\bigl(x_2-x_1^2x_2\bigr)^2+\bigl(-x_1+x_1x_2^2\bigr)^2-x_1^2-x_2^2\\
&=
x_1^2x_2^2\bigl(x_1^2+x_2^2-4\bigr).
\end{align*}
Quindi, in un intorno sufficientemente piccolo dell'origine,
\[
\Delta V \le 0.
\]
Per esempio, nel disco
\[
x_1^2+x_2^2<4
\]
si ha effettivamente \(\Delta V\le 0\).

Tuttavia
\[
\Delta V=0
\]
non si verifica solo nell'origine, ma anche sugli assi \(x_1=0\) e \(x_2=0\). Bisogna quindi studiare il massimo insieme invariante contenuto in tale insieme.

Se partiamo da un punto dell'asse \(x_1\), per esempio
\[
x(0)=(a,0),
\]
si ottiene:
\[
x(1)=(0,-a),
\qquad
x(2)=(-a,0),
\qquad
x(3)=(0,a),
\qquad
x(4)=(a,0).
\]
Quindi compare un'orbita periodica di periodo \(4\).

Dunque il massimo insieme invariante contenuto in \(\{\Delta V=0\}\) non è il solo punto \((0,0)\).  
Perciò LaSalle non dà asintotica stabilità. In effetti l'origine risulta soltanto \textbf{stabile}, non asintoticamente stabile.

\medskip

\noindent
Questo esempio è importante perché mostra una differenza sostanziale tra:
\begin{itemize}
    \item ``\(\Delta V\le 0\)'',
    \item ``\(\Delta V<0\)''.
\end{itemize}
Nel primo caso, senza studiare l'insieme invariante, non si può concludere la convergenza all'equilibrio.

\section{Lyapunov e sintesi del controllo}

Una delle idee più importanti della lezione è che i metodi di Lyapunov non servono solo per \emph{analizzare} la stabilità, ma anche per \emph{progettare} controlli stabilizzanti.

\subsection{Idea di fondo}

Supponiamo di avere un sistema non lineare controllato:
\[
\dot x = f(x,u).
\]
Se vogliamo stabilizzare un equilibrio \((\bar x,\bar u)\), conviene traslare il sistema:
\[
z=x-\bar x,
\qquad
v=u-\bar u.
\]
In queste coordinate l'equilibrio da stabilizzare diventa l'origine:
\[
z=0,\qquad v=0.
\]

L'idea è allora cercare una funzione \(V(z)\) positiva definita e scegliere una legge di controllo
\[
v=\alpha(z)
\]
in modo che
\[
L_fV(z)\le 0
\quad \text{o meglio} \quad
L_fV(z)<0.
\]

In questo modo il controllo è \emph{stabilizzante per costruzione}.

\subsection{Perché il controllo non deve spostare l'equilibrio}

Se il controllo serve a stabilizzare l'equilibrio, è essenziale che nel punto di equilibrio esso si annulli in coordinate traslate:
\[
v(0)=0.
\]
Equivalentemente, in coordinate originali deve valere
\[
u(\bar x)=\bar u.
\]
Se così non fosse, il controllo introdurrebbe una forza costante anche nel punto che vogliamo mantenere come equilibrio, spostandolo.

\section{Controllare prima e linearizzare dopo, o viceversa}

Per un sistema non lineare
\[
\dot x=f(x,u),
\]
sia \((\bar x,\bar u)\) un equilibrio, cioè
\[
f(\bar x,\bar u)=0.
\]
Linearizzando attorno a \((\bar x,\bar u)\), otteniamo
\[
\dot z = Az+Bv,
\]
dove
\[
A=\left.\frac{\partial f}{\partial x}\right|_{(\bar x,\bar u)},
\qquad
B=\left.\frac{\partial f}{\partial u}\right|_{(\bar x,\bar u)}.
\]

Se imponiamo una retroazione lineare locale
\[
v=Kz,
\]
il linearizzato in anello chiuso diventa
\[
\dot z = (A+BK)z.
\]

Questa osservazione è fondamentale: localmente, progettare il feedback sul sistema linearizzato equivale a linearizzare il sistema già controllato.  
Perciò, se \(A+BK\) è Hurwitz, il sistema non lineare controllato è localmente asintoticamente stabile attorno all'equilibrio, per il metodo indiretto di Lyapunov.

\section{Esempio di sintesi con Lyapunov}

Consideriamo il sistema
\begin{equation}
\label{eq:L23-controllo}
\begin{cases}
\dot x_1 = (x_2+1)^3,\\[4pt]
\dot x_2 = (x_1-1)^3 + u + 2.
\end{cases}
\end{equation}

Vogliamo scegliere \(u(t)\) in modo da rendere stabile l'equilibrio
\[
(\bar x_1,\bar x_2)=(1,-1).
\]

\subsection{Traslazione}

Poniamo
\[
z_1=x_1-1,\qquad z_2=x_2+1,
\qquad
v=u+2.
\]
In questo modo l'equilibrio desiderato diventa l'origine:
\[
z=0,\qquad v=0.
\]
Il sistema si riscrive come
\[
\begin{cases}
\dot z_1 = z_2^3,\\[4pt]
\dot z_2 = z_1^3 + v.
\end{cases}
\]

\subsection{Scelta della funzione di Lyapunov}

Proviamo con
\[
V(z)=\frac12(z_1^2+z_2^2).
\]
Allora
\[
L_fV=z_1\dot z_1+z_2\dot z_2
= z_1 z_2^3 + z_2(z_1^3+v).
\]
Sviluppando:
\[
L_fV = z_1z_2^3 + z_1^3z_2 + z_2v.
\]
I primi due termini hanno segno indefinito, quindi da soli non consentono di concludere nulla.  
L'idea è usare il controllo per cancellarli e aggiungere un termine dissipativo.

Scegliamo
\[
v = -z_1z_2^2 - z_1^3 - z_2.
\]
Allora
\begin{align*}
L_fV
&= z_1z_2^3 + z_1^3z_2 + z_2\bigl(-z_1z_2^2-z_1^3-z_2\bigr)\\
&= -z_2^2 \le 0.
\end{align*}
Quindi il sistema controllato è stabile.

Per vedere l'asintotica stabilità, osserviamo che
\[
L_fV=0 \iff z_2=0.
\]
Se \(z_2=0\), allora
\[
\dot z_1=0,
\qquad
\dot z_2=z_1^3.
\]
Affinché la traiettoria resti in \(\{z_2=0\}\), deve essere
\[
z_1^3=0 \quad \Longrightarrow \quad z_1=0.
\]
Quindi il massimo insieme invariante contenuto in \(\{L_fV=0\}\) è il solo punto
\[
M=\{0\}.
\]
Per LaSalle, l'origine del sistema controllato è asintoticamente stabile.

\subsection{Legge di controllo in coordinate originali}

Ricordando che
\[
z_1=x_1-1,\qquad z_2=x_2+1,\qquad v=u+2,
\]
si ottiene
\[
u = v-2,
\]
quindi una legge di controllo stabilizzante è
\[
\boxed{
u(x)= -2 - (x_1-1)(x_2+1)^2 - (x_1-1)^3 - (x_2+1).
}
\]

\section{Cicli limite: cenno qualitativo}

Nella parte finale della lezione compare un'osservazione importante sui sistemi non lineari: oltre agli equilibri, possono esistere anche \textbf{cicli limite}.

\medskip

\noindent
\textbf{Definizione informale.}
Un ciclo limite è una traiettoria periodica isolata nello spazio delle fasi.  
In altre parole, il sistema torna periodicamente sugli stessi stati e la traiettoria descrive una curva chiusa.

\medskip

\noindent
Questo fenomeno è tipico dei sistemi non lineari.  
Nei sistemi lineari autonomi non banali, infatti, le traiettorie periodiche non sono isolate nello stesso senso: se esiste una traiettoria chiusa, tipicamente ne esiste un'intera famiglia.

\medskip

\noindent
Un ciclo limite può essere:
\begin{itemize}
    \item attrattivo,
    \item repulsivo,
    \item semistabile.
\end{itemize}

Talvolta, per studiarne l'attrattività, si costruisce una funzione del tipo
\[
V(x)=C(x)^2,
\]
dove la curva
\[
C(x)=0
\]
descrive il ciclo o una curva invariante candidata.  
L'idea è che \(V\) sia nulla sulla curva e positiva fuori, così da misurare la distanza energetica dalla traiettoria chiusa.

\section{Schema operativo per gli esercizi di fine lezione}

Gli esercizi di questa lezione si risolvono in genere seguendo questo schema.

\subsection*{A. Trovare gli equilibri}
Si impone:
\[
f(x)=0
\]
oppure, nel caso controllato,
\[
f(x,u)=0
\]
con l'ingresso assegnato o da scegliere.

\subsection*{B. Provare il linearizzato}
Si calcola il Jacobiano:
\[
A=\frac{\partial f}{\partial x}.
\]
Se tutti gli autovalori hanno parte reale negativa (o modulo minore di \(1\) nel discreto), allora si conclude subito la stabilità asintotica locale.

\subsection*{C. Se il linearizzato è inconcludente, usare Lyapunov}
Si prova una funzione:
\begin{itemize}
    \item quadratica, se il sistema lo suggerisce;
    \item energetica, se c'è una struttura meccanica;
    \item con potenze adattate alla non linearità, se compaiono monomi di grado alto.
\end{itemize}

\subsection*{D. Se \(L_fV\) o \(\Delta V\) è solo semidefinita, usare LaSalle}
Si calcola
\[
E=\{x:L_fV(x)=0\}
\qquad \text{oppure} \qquad
E=\{x:\Delta V(x)=0\},
\]
poi si cerca il massimo insieme invariante contenuto in \(E\).

\subsection*{E. Nei problemi di controllo, traslare prima l'equilibrio}
Se l'equilibrio desiderato non è l'origine:
\[
z=x-\bar x,\qquad v=u-\bar u.
\]
Poi si progetta il controllo in \((z,v)\), imponendo
\[
v(0)=0.
\]

\section{Riassunto finale}

La lezione porta a quattro idee chiave.

\begin{enumerate}
    \item \textbf{LaSalle completa Lyapunov:} se \(V\) non cresce, la traiettoria converge al massimo insieme invariante contenuto in \(\{L_fV=0\}\).
    \item \textbf{Per concludere l'asintotica stabilità} basta spesso mostrare che tale massimo insieme invariante è il solo equilibrio.
    \item \textbf{Le funzioni di Lyapunov si possono progettare}, non solo usare a posteriori: questo apre direttamente alla sintesi del controllo.
    \item \textbf{Nei sistemi non lineari} compaiono fenomeni più ricchi degli equilibri, come i cicli limite, che richiedono strumenti qualitativi più raffinati.
\end{enumerate}

\medskip

\noindent
In forma compatta, il messaggio da ricordare è:

\begin{center}
\emph{Se trovi una funzione \(V\) che non aumenta, non hai ancora finito: devi chiederti dove il sistema può restare quando \(V\) smette di diminuire. La risposta la dà il massimo insieme invariante di LaSalle.}
\end{center}